\sectionkicker{Source-backed technical notes}
\subsection*{Multimodalは生成からAction Spaceへ — Video・World Model・Robotics: Technical Notes}
本欄は記事本文で圧縮した一次資料上の情報を、比較・再検証しやすい形で整理したものである。時系列、確認済みの主張、留意点、一次資料URLを掲載する。
\medskip
\subsection*{Theme at a glance}
\addcontentsline{toc}{subsection}{Theme at a glance}
\begin{center}
\footnotesize
\begin{tabularx}{\linewidth}{@{}p{0.20\linewidth}p{0.10\linewidth}p{0.14\linewidth}X@{}}
\toprule
資料 & 位置づけ & 種別 & 時系列 \\
\midrule
Google DeepMind H2 world-model, robotics, and agent research & 主要資料 & モデル更新 & 2025-08-05 (Genie 3（発表）); 2025-09-25 (Gemini Robotics 1.5（発表）); 2025-11-13 (SIMA 2（発表）); 2025-11-18 (Gemini 3（発表）)      \\
Sora 2 & 補足資料 & モデル & 2025-09-30 (映像・音声モデル（公開）); 2025-09-30 (System Card公開)      \\
MiniMax M2 / Hailuo 2.3 / M2.1 lifecycle & 補足資料 & モデル更新 & 2025-10-27 (MiniMax M2（公開）); 2025-10-28 (Hailuo 2.3（公開）); 2025-12-22 (MiniMax M2.1（公開）)      \\
\bottomrule
\end{tabularx}
\end{center}
\normalsize

\begin{technicalnote}{Google DeepMind H2 world-model, robotics, and agent research}{主要資料}
\begingroup
% reader-facing Technical Notes break policy
\widowpenalty=10000
\clubpenalty=10000
\displaywidowpenalty=10000
\begin{tabularx}{\linewidth}{@{}>{\bfseries}p{0.22\linewidth}X@{}}
組織 & Google DeepMind \\
種別 & モデル更新 \\
時系列 & 2025-08-05 (Genie 3（発表）); 2025-09-25 (Gemini Robotics 1.5（発表）); 2025-11-13 (SIMA 2（発表）); 2025-11-18 (Gemini 3（発表）) \\
\end{tabularx}
\smallskip
{\bfseries 一次資料から整理したtechnical points}
\begin{itemize}[leftmargin=1.5em,itemsep=0.35em]
\item \textbf{一次情報で確認できる事実}: Google DeepMindは8月5日、リアルタイムで対話可能なworld modelとしてGenie 3を発表した。
\item \textbf{一次情報で確認できる事実}: Google DeepMindは9月25日、物理行動と身体化推論・ツール利用を組み合わせるGemini Robotics 1.5とRobotics-ER 1.5を発表した。Robotics-ER 1.5はGemini APIで提供され、Robotics 1.5は一部パートナー向けに限定された。
\item \textbf{一次情報で確認できる事実}: 11月13日のSIMA 2はGeminiの推論を3D世界のエージェントへ統合し、11月18日にはより広範なfrontier modelとしてGemini 3が続いた。
\end{itemize}
\begin{minipage}{\linewidth}
% reader-facing Technical Notes generic boundary/source tail group
{\bfseries 読む際の境界}
\begin{itemize}[leftmargin=1.5em,itemsep=0.35em]
\item \textbf{分析上の留意点}: 保存したDeepMindのindexスナップショットはページ分割されており、先頭ページだけでは2025年後半の完全な記録にならないため、個別の公式ページを用いて検証した。
\item \textbf{分析上の留意点}: 研究プロトタイプ、一部パートナー限定のロボティクス提供、一般提供されるAPIモデルは、それぞれ異なる提供状態として区別する。
\end{itemize}
\begin{samepage}
{\bfseries 一次資料}
\begingroup\sloppy
\Urlmuskip=0mu plus 2mu\relax
\begin{itemize}[leftmargin=1.5em,itemsep=0.25em]
\item {\scriptsize\url{https://deepmind.google/blog/}}
\end{itemize}
\endgroup
\end{samepage}
\end{minipage}
\endgroup
\end{technicalnote}
\medskip

\begin{technicalnote}{Sora 2}{補足資料}
\begingroup
% reader-facing Technical Notes break policy
\widowpenalty=10000
\clubpenalty=10000
\displaywidowpenalty=10000
\begin{tabularx}{\linewidth}{@{}>{\bfseries}p{0.22\linewidth}X@{}}
組織 & OpenAI \\
種別 & モデル \\
時系列 & 2025-09-30 (映像・音声モデル（公開）); 2025-09-30 (System Card公開) \\
\end{tabularx}
\smallskip
{\bfseries 一次資料から整理したtechnical points}
\begin{itemize}[leftmargin=1.5em,itemsep=0.35em]
\item \textbf{一次情報で確認できる事実}: OpenAIは9月30日にSora 2を公開し、同日にSora 2 System Cardを公開した。
\item \textbf{一次情報で確認できる事実}: release packageでは映像生成と音声を同じモデル／製品更新の一部として扱い、System Cardはその提供時点の評価・control境界を別資料として示した。
\end{itemize}
\begin{minipage}{\linewidth}
% reader-facing Technical Notes generic boundary/source tail group
{\bfseries 読む際の境界}
\begin{itemize}[leftmargin=1.5em,itemsep=0.35em]
\item \textbf{分析上の留意点}: 「Sora 2」の扱いでは、一次資料で確認できる範囲、提供時点、評価条件を維持し、記載範囲を超えて一般化しない。
\end{itemize}
\begin{samepage}
{\bfseries 一次資料}
\begingroup\sloppy
\Urlmuskip=0mu plus 2mu\relax
\begin{itemize}[leftmargin=1.5em,itemsep=0.25em]
\item {\scriptsize\url{https://openai.com/index/sora-2}}
\item {\scriptsize\url{https://openai.com/index/sora-2-system-card}}
\end{itemize}
\endgroup
\end{samepage}
\end{minipage}
\endgroup
\end{technicalnote}
\medskip

\begin{technicalnote}{MiniMax M2 / Hailuo 2.3 / M2.1 lifecycle}{補足資料}
\begingroup
% reader-facing Technical Notes break policy
\widowpenalty=10000
\clubpenalty=10000
\displaywidowpenalty=10000
\begin{tabularx}{\linewidth}{@{}>{\bfseries}p{0.22\linewidth}X@{}}
組織 & MiniMax \\
種別 & モデル更新 \\
時系列 & 2025-10-27 (MiniMax M2（公開）); 2025-10-28 (Hailuo 2.3（公開）); 2025-12-22 (MiniMax M2.1（公開）) \\
\end{tabularx}
\smallskip
{\bfseries 一次資料から整理したtechnical points}
\begin{itemize}[leftmargin=1.5em,itemsep=0.35em]
\item \textbf{一次情報で確認できる事実}: MiniMaxの公式リリースノートには、10月27日のエージェント指向モデルMiniMax M2、10月28日の動画モデルHailuo 2.3／2.3-Fast、12月22日のMiniMax M2.1が記録されている。
\item \textbf{一次情報で確認できる事実}: Hailuo 2.3はHailuo製品とOpen Platform APIの双方で展開され、モデル更新とMedia Agentワークフローが組み合わされた。
\end{itemize}
\begin{minipage}{\linewidth}
% reader-facing Technical Notes generic boundary/source tail group
{\bfseries 読む際の境界}
\begin{itemize}[leftmargin=1.5em,itemsep=0.35em]
\item \textbf{分析上の留意点}: 保存したminimax.io/newsのランディングページはクライアント描画で情報が疎なため、検証には個別のMiniMax公式リリースノート／ニュースを用いた。
\end{itemize}
\begin{samepage}
{\bfseries 一次資料}
\begingroup\sloppy
\Urlmuskip=0mu plus 2mu\relax
\begin{itemize}[leftmargin=1.5em,itemsep=0.25em]
\item {\scriptsize\url{https://www.minimax.io/news}}
\end{itemize}
\endgroup
\end{samepage}
\end{minipage}
\endgroup
\end{technicalnote}
\medskip
